\documentclass{ximera}




\title{A Basic Review of Derivatives}   % Use xmlatex all -s to compile when SERVE refuses to work.
\author{Kelly Stady}
\license{CC: 0}         % replace with an appropriate license, or set it in xmPreamble

\begin{document}
\begin{abstract}
    A basic review of derivatives for students who have practiced finding derivatives using the limit definition, basic differentiation rules, product rule, quotient rule
    and chain rule.
\end{abstract}
\maketitle
\label{xim:DerivativeActivity}



\section*{The Limit Definition of the Derivative}

Recall that the slope of a curve at a point is equal to the slope of the tangent line at that point.  The figure below shows the graph
of a function $f(x)$ and its tangent line at point $P (x, f(x)).$  To find the slope of the tangent line at $P$ algebraically, we would need to know two points that
lie on the line.  The tangent line and curve have one point in common, but we don't know of any other points that lie on the tangent line.
So, instead of finding the exact slope we will approximate it using the slope of a secant line.  Recall that a \textbf{secant line} is a line that passes
through two points on a curve.

The slope of $f$ at point $P(x, f(x))$ can be approximated by drawing a secant line through $P$ and
a nearby point $Q(x+\Delta x, f(x+\Delta x)).$  If $\Delta x$ is a small number,
then the slope of the secant line will provide a decent approximation to the slope of the
tangent line.
\[ m_{\tan} \approx m_{\sec} \]


\begin{center}
\begin{tikzpicture}
\begin{axis}[
    axis line style = {thick},
    axis lines = middle,
    xmin = -1.0, xmax = 3.5,
    ymin = -2.0, ymax = 15,
    % Use \strut to align the baselines of the labels perfectly
    xtick = {1.2, 2.2},
    xticklabels = {\strut $x$, \strut $x + \Delta x$},
    ytick = {3.32, 9.02},
    yticklabels = {$f(x)$, $f(x + \Delta x)$},
    xlabel = {$x$}, ylabel = {$y$},
    every axis x label/.style={at={(current axis.right of origin)}, anchor=west},
    every axis y label/.style={at={(current axis.above origin)}, anchor=south},
    samples = 100,
    clip = false
]
    % 1. The Curve
    \addplot [black, very thick, domain=-0.8:2.7] {exp(x)};

    % 2. Tangent Line: Blue
    \addplot [glaucous!90!black, thick, domain=-0.2:3.2] {3.3201*x - 0.6640}
        node[pos=1, right, yshift=5pt, xshift=-4.3pt] {\small \textbf{Tangent}};

    % 3. Secant Line: Orange
    \addplot [orange!80!black, dashed, thick, domain=0.5:3.05] {5.70*x - 3.52}
        node[pos=1, right, yshift=5.3pt, xshift=-1.2pt] {\small \textbf{Secant}};

    % 4. Points P and Q
    \addplot[only marks, mark=*, mark size=1.8pt] coordinates {(1.2, 3.32) (2.2, 9.02)};
    \node[anchor=south, yshift=3pt] at (1.2, 3.32) {$P$};
    \node[anchor=south, yshift=3pt] at (2.15, 9.02) {$Q$};

    % 5. Dotted lines to axes
    \draw [gray, thick, dotted] (1.2, 0) -- (1.2, 3.32) -- (0, 3.32);
    \draw [gray, thick, dotted] (2.2, 0) -- (2.2, 9.02) -- (0, 9.02);

\end{axis}
\end{tikzpicture}

% Alt Text
\xmalt{A graph shows a black exponential curve rising from left to right. Two points are marked on the curve: point P at $x$ and
point Q further to the right at $x + \Delta x.$ A solid blue line labeled "Tangent" touches the curve exactly at point P. A dashed
orange line labeled "Secant" passes through both points P and Q. Gray dotted lines drop from both points to the $x$ and $y$ axes to
mark their coordinates as $x, x + \Delta x, f(x)$ and $f(x + \Delta x).$}
\end{center}

Notice that as $\Delta x \rightarrow 0,$ point $Q$ will get closer to point $P,$ and the secant line will start to look more like
the tangent line.  Therefore, the slope of the tangent line is the limit of the slope of the secant line.
\[ m_{\sec} = \frac{f(x+\Delta x) - f(x)}{\Delta x}  \hspace{0.3in}  \textbf{Slope of the Secant Line} \]

\[ m_{\tan} = \lim_{\Delta x \rightarrow 0} \frac{f(x+\Delta x) - f(x)}{\Delta x} \hspace{0.2in}  \textbf{Slope of the Tangent Line} \]


Below is an interactive Desmos graph.  Click and drag point $Q$ toward point $P$ to see what happens when $\Delta x \rightarrow 0.$
If the left sidebar is showing the math, click $\mathbf{\ll}$ to hide the contents.


\begin{center}
\desmos{zsdgpux7ia}{800}{600}              % \desmos{ID}{width}{height}  dimensions are in pixels

\xmalt{An interactive Desmos graph shows the function $f(x) = e^x.$ A fixed point $P$ is marked in quadrant I and a movable point $Q$ is
positioned to its right at $x + \Delta x.$ There is a solid blue tangent line at $P$ and a dashed orange secant line passing through both $P$
and $Q.$ As the user drags point $Q$ toward $P,$ $\Delta x \rightarrow 0$ and the orange secant line approaches the blue tangent line.}

\end{center}

\end{document}

